 Un vídeo digital se define por un conjunto de propiedades objetivas que determinan tanto su representación técnica como la naturaleza de su contenido. Estas propiedades, denominadas habitualmente características del vídeo, permiten describir, comparar y procesar cualquier secuencia de manera cuantitativa, con independencia del codificador o de la plataforma de reproducción empleada.

 Del conjunto de características que componen el vídeo, pueden diferenciarse dos grupos. El primero agrupa los parámetros de formato, que definen cómo se muestrea y se representa la señal, siendo los más relevantes la resolución (medida en píxeles), la frecuencia de fotogramas, la tasa de bits, la profundidad de color (número de bits por muestra), el rango dinámico (SDR o HDR) y la relación de aspecto. Estos parámetros son configurables y determinan el volumen de datos, los requisitos de cómputo y el ancho de banda necesarios para la transmisión.

 En el segundo grupo se encuentran las características de contenido, que describen la complejidad intrínseca de la escena y no dependen directamente de los parámetros de formato, sino de lo que aparece en cada imagen y de cómo evoluciona a lo largo del tiempo. Dentro de este grupo destacan dos dimensiones complementarias: la complejidad espacial (SI), asociada a la cantidad de detalle, textura y bordes presentes en cada fotograma; y la complejidad temporal (TI), asociada a la magnitud del movimiento y a los cambios entre fotogramas consecutivos. Estas características son las que condicionan el esfuerzo real del codificador, ya que determinan el número de candidatos de predicción, la dificultad de la estimación de movimiento y el volumen de residuos a procesar.

 \subsection{Spatial Information y Temporal Information}

 En la actualidad, las métricas SI y TI representan el estado del arte en la caracterización del vídeo y se emplean para determinar su complejidad. Ambos parámetros se encuentran estandarizados por la UIT-T bajo la Recomendación \textbf{P.910: Subjective video quality assessment methods for multimedia applications} \cite{itu-quality-metrics}, cuya versión vigente data de 10/2023, a la espera de la publicación oficial de la nueva versión fechada en 07/2026. Cabe destacar que, a diferencia de las versiones previas (que empleaban el valor máximo como medida de agregación), la versión actual recomienda utilizar la media.

 Antes de calcular los valores de SI y TI, a cada fotograma del vídeo se le aplica un conjunto de transformaciones:
 \begin{enumerate}
   \item Extracción de la luminancia (ignorando la información de color)
   \item Normalización a valores en el rango \([0,1]\).
   \item Escalado de rango escalado para ajustar y expandir los valores de luminancia de secuencias de vídeo con rango de color limitado al rango completo (si el vídeo tiene un rango de color completo este paso no se realiza).
   \item A continuación, se aplica la función de transferencia electro-óptica (EOTF) para convertir los valores de luma en valores de luminancia dentro del dominio físico. Esto permite traducir las señales digitales a la cantidad de luz real emitida por una pantalla de referencia según el rango dinámico del vídeo.
 \end{enumerate}
 
 A partir de esta información, se obtiene un valor de SI para cada fotograma aplicando un filtro de Sobel (detección de bordes) y se calcula la desviación estándar. Finalmente, el valor de SI de la secuencia se obtiene como la media de los valores obtenidos para cada fotograma.
 
 El valor de TI para cada fotograma se obtiene aplicando la diferencia entre el fotograma $i$ y el fotograma $i-1$ y se calcula la desviación estándar. El valor de TI, igual que en el caso de SI, se obtiene como la media de los valores de TI obtenidos para cada fotograma.

 Para finalizar, es necesario destacar que el cálculo de las métricas SI y TI resulta costoso a nivel temporal y energético, según los resultados hallados y expuestos en el artículo \textbf{Green Video Complexity Analysis for Efficient Encoding in Adaptive Video Streaming} \cite{video-complexity}. En dicho artículo se propone un marco alternativo que permite calcular el valor de las métricas Spatial Information y Temporal Information con un consumo energético y un coste temporal sustancialmente menor.

 En relación con este trabajo de investigación, se pretende comprobar si existe una correlación entre el SI/TI de cada vídeo del dataset y el consumo energético. Además, se analizarán los resultados de consumo de cada proceso de codificación con cada uno de los códecs empleados, con el fin de detectar la influencia del SI y/o del TI sobre los códecs en caso de que dicha influencia exista.
